\section{Introduction}
Receipt Key Information Extraction (KIE) — recovering structured
fields such as merchant name, date, address, and total from scanned
receipts — underpins downstream financial automation. Two paradigms
have emerged: (i) end-to-end vision-language models that treat KIE as
image-conditioned text generation~\cite{kim2022donut}, and (ii)
detect-then-read pipelines that first locate text regions with an
object detector~\cite{jocher2023yolov8} and then transcribe each
region with a specialised OCR model~\cite{li2023trocr}.

The central question this paper asks is: \emph{does the added
modularity of a pipeline justify its lower end-to-end F1, and can a
learned field assigner close the gap?} We answer on the SROIE
dataset~\cite{huang2019icdar}, which provides 626 labelled receipt
images with four ground-truth fields per image.
